\section{Conclusion}\label{sec:conclusion}
Neural Bellman Operators connect a specified continuation object with feasible
economic decisions. The nonlinear capital economy makes this connection
explicit: an exact finite Bellman identity accounts for continuation error
along the chosen action change, while a separate diffusion account assesses
the deployed continuous-time policy. Repeated decisions permit the same
continuation representation to be used across economic tasks. A fixed
gradient-dictionary calculation identifies the approximation--variance
tradeoff of a linear projection, and the numerical experiment measures the
economic performance of the nonlinear fitted procedure.

The complete continuation-menu experiment establishes material NBO payoff
gains over materialized Raw and direct-policy actors in two previously
unpiloted fifty-dimensional economies. NBO also uses less construction and
query work than direct policy in those cells. At the final prescribed SAA
budget, six of eight economic cells have equivalent payoffs within
$10^{-4}$ and lower construction and query work for NBO. Cached SAA has
materially higher payoffs in both longest-horizon cells. The full set of
earlier budgets, cheaper vector and Raw actors, and common confirmation
costs remains part of the evidence. This result identifies specific
economic gains from a reusable fitted continuation and the work needed
to obtain them.

Independent confirmation supports the original-policy paired comparison
under a prospectively fixed analysis. The strengthened HJB procedure shows
why the implementation must remain explicit: NBO's material advantage
over both new HJB deployments is established in the low-volatility,
fifty-dimensional economy, while the remaining direct comparisons are
unresolved. A new signed continuation statistic verifies positive
occupation Bellman gains by retaining the common-path cancellation.
These positive payoff statements coexist with the reported full-class
regret bounds and the finite menu's separate scalar accuracy account.

The continuation-risk comparison supplies separate evidence about the fitted nonlinear object. Cancellation of common simulation noise makes eight risk advantages and four reversals identifiable across the complete 24-event family. The economic payoff experiments, rather than risk ordering alone, establish the value of the returned decisions. Exact closure of the complete menu intervals further identifies the least-cost certified candidate at the common accuracy tolerance. NBO attains that distinction in the fifty-dimensional Intermediate design; other designs select simulation or cheaper fitted actors. The result concerns an explicit catalogue of computed procedures and observed accounting costs. It does not replace the broader adapted-class regret account or the full bill for producing the comparison.

The economic applications retain their own utility domains, boundary
conditions, and deviation criteria. These determine the continuation
that must be learned and the decision that constitutes an improvement.
The current application companion preserves the complete models,
derivations, proofs, and numerical records for recursive utility,
endogenous preferences, temporal selves, and strategic interaction.
The common evaluation--improvement structure makes their computational
requirements comparable while retaining their distinct economic content.
